problem:  
Let \(M^n \subset S^{n+p}(1)\) be a compact minimal submanifold with second fundamental form \(B = \{h_{ij}^\alpha\}\), and denote its shape operators by \(A_\alpha\) in each normal direction \(e_\alpha\).  
The **Simons identity** contains a commutator term \(\sum_{\alpha,\beta}\|[A_\alpha,A_\beta]\|^2\), which vanishes exactly when the family \(\{A_\alpha\}\) of shape operators is simultaneously diagonalizable—equivalently, when \(M\) is *isoparametric in the normal directions*.  For known homogeneous minimal examples (such as Clifford and Veronese-type immersions), this commutator term is constant on \(M\).

Define the **commutator energy functional**
\[
\mathcal{C}(M) = \int_M \sum_{\alpha,\beta=1}^p \|[A_\alpha,A_\beta]\|^2\, d\mathrm{vol}_M.
\]
It measures the total failure of the shape operators to commute.

**Problem (Commutator–Curvature Gap Conjecture):**  
For each pair \((n,p)\), do there exist positive constants \(\varepsilon_{n,p},\,\delta_{n,p}>0\) such that if  
\[
M^n \subset S^{n+p}(1)
\]
is a compact minimal submanifold satisfying simultaneously
\[
\|B\|^2 < n + \varepsilon_{n,p}, \qquad \mathcal{C}(M) < \delta_{n,p},
\]
then \(M\) must be congruent (up to an isometry of \(S^{n+p}(1)\)) to one of the known **isoparametric minimal submanifolds**, i.e. those for which all \(A_\alpha\) commute and have constant spectral type?

Equivalently: is there a *universal commutator gap* around the Simons curvature pinching boundary in which every minimal submanifold whose shape operators are “almost commuting’’ must in fact be globally isoparametric (and hence homogeneous)?

Why is it a "good" problem:

1. **Geometrically natural and mathematically precise:**  
   The functional \(\mathcal{C}(M)\) arises directly from the Simons identity, not from any ad hoc construction.  It quantifies a deep geometric property—how far the family of shape operators is from being simultaneously diagonalizable—while \(\|B\|^2\) measures total curvature.  Both are intrinsic to the established analytical framework.

2. **Conceptual simplicity yet deep implications:**  
   The conjecture is stated in ordinary geometric language—“small total curvature plus almost commuting shape operators implies exact homogeneity”—but its resolution would require mastery of the full interplay between the submanifold’s extrinsic algebra and curvature analysis.

3. **Bridges multiple subfields:**  
   The problem unites methods from geometric analysis (Simons-type Laplacian identities), algebraic geometry and representation theory (classification of isoparametric submanifolds), and global differential geometry (curvature pinching).  Any advance would necessarily entwine these perspectives.

4. **Close continuity with known results yet genuinely open:**  
   The known Simons gap theorems treat scalar curvature bounds; no theorem yet quantifies a *commutator gap*.  Discovering such a relation would refine the concept of rigidity from magnitude to algebraic structure.

5. **Potential transformative consequence:**  
   A positive answer would show that, infinitesimally, *nearly commuting shape operators force full symmetry*, guiding future classification of minimal submanifolds in spheres.  A counterexample would reveal unexpected families violating such rigidity, reshaping our understanding of geometric pinching.

This problem is “good’’ because it is crisp, grounded in intrinsic equations of minimal geometry, and pushes the frontier of our understanding of algebraic and analytic rigidity in positively curved ambient spaces.